\subsubsection*{3.2.5 Final Design Class Diagram}
\addcontentsline{toc}{subsubsection}{3.2.5 Final Design Class Diagram}
\pdfbookmark[2]{3.2.5 Final Design Class Diagram}{p7h80}
Author: 722I0004 \ensuremath{\cdot} AI-assisted (Claude)

The final design class diagram refines the layer views of 3.2.2 with what the sequence diagrams in 3.2.4 and the implementation settled. Figure 3.2.5A is an overview of all classes; the detail sheets after it give the members of one controller at a time.

\begin{itemize}
\item \textbf{Overview first, then detail sheets.} Figure 3.2.5A shows every class of the four layers by name, with the calls, the \guillemotleft{}persists\guillemotright{} and \guillemotleft{}loads\guillemotright{} links and the domain associations, drawn as plain lines; their multiplicities and navigability are in the domain views of 3.2.2 and the domain model (2.1.7), because the labels pile up at the small name boxes. The detail sheets leave the associations out. The members are too many to print readably on the same page, so the detail sheets take one controller at a time and show the view functions that call it and the DAOs and domain classes it uses, with member signatures. A controller that needs more than one sheet has them numbered (A1, A2 and so on); each caption says what the sheet shows. Each figure is laid out for its page so that the estimated smallest text is 9 pt or more.

\item \textbf{Only the members this use case calls.} The sequence diagrams fixed which operations each class must offer; the complete member lists stay in the 3.2.1 tables and 3.2.2 views.

\item \textbf{Signatures read from the code.} Parameter names and types are read from the source by \texttt{models/\allowbreak{}shared/\allowbreak{}final\_\allowbreak{}dcd.\allowbreak{}py} when the figures are generated, and each member is matched with the operation rows of the design model. This is a cross-check at generation time, not a guarantee: the figures must be regenerated after a code change.

\item \textbf{Navigability and persistence.} Layers run top to bottom (view, controller, data access, domain) and the dashed arrows run the same way: view to controller, controller to the DAOs and domain classes it uses, and DAO to the entity it maps. \guillemotleft{}persists\guillemotright{} marks a DAO that writes the entity's table in this use case (an INSERT, UPDATE or DELETE in a method the use case calls); \guillemotleft{}loads\guillemotright{} marks one that only reads it. The arrow points from the DAO to the entity, so the domain layer does not depend on data access. Domain associations keep their multiplicities from the domain model (2.1.7).

\end{itemize}
\begin{landscape}
\begin{figure}[p]
\centering
\includegraphics[width=230mm,height=150mm,keepaspectratio]{\detokenize{models/UC-02/dcd_uc02_final_overview.png}}
\caption{Figure 3.2.5A. UC-02 final design class diagram, overview (class names only) --- Author: 722I0004 \ensuremath{\cdot} AI-assisted (Claude).}\label{P7-0615}
\end{figure}
\end{landscape}
\begin{landscape}
\begin{figure}[p]
\centering
\includegraphics[width=230mm,height=150mm,keepaspectratio]{\detokenize{models/UC-02/dcd_uc02_final_detail_A1.png}}
\caption{Figure 3.2.5B. UC-02 final design class diagram, detail sheet A1 (DeliveryController, view and controller) --- Author: 722I0004 \ensuremath{\cdot} AI-assisted (Claude).}\label{P7-0641}
\end{figure}
\end{landscape}
\begin{figure}[p]
\centering
\includegraphics[width=170mm,height=225mm,keepaspectratio]{\detokenize{models/UC-02/dcd_uc02_final_detail_A2.png}}
\caption{Figure 3.2.5C. UC-02 final design class diagram, detail sheet A2 (DeliveryController, controller and domain) --- Author: 722I0004 \ensuremath{\cdot} AI-assisted (Claude).}\label{P7-0642}
\end{figure}
\begin{figure}[p]
\centering
\includegraphics[width=170mm,height=225mm,keepaspectratio]{\detokenize{models/UC-02/dcd_uc02_final_detail_A3.png}}
\caption{Figure 3.2.5D. UC-02 final design class diagram, detail sheet A3 (DeliveryController, data access, part 1 of 2) --- Author: 722I0004 \ensuremath{\cdot} AI-assisted (Claude).}\label{P7-0643}
\end{figure}
\begin{landscape}
\begin{figure}[p]
\centering
\includegraphics[width=230mm,height=150mm,keepaspectratio]{\detokenize{models/UC-02/dcd_uc02_final_detail_A4.png}}
\caption{Figure 3.2.5E. UC-02 final design class diagram, detail sheet A4 (DeliveryController, data access, part 2 of 2) --- Author: 722I0004 \ensuremath{\cdot} AI-assisted (Claude).}\label{P7-0644}
\end{figure}
\end{landscape}
Behaviour refined between the first design and the final design (each refinement is covered by the tests in 5.1). Every state-changing operation runs in one \texttt{BEGIN IMMEDIATE} transaction and checks the row count of each guarded UPDATE, so a stale page, a guessed ID or a competing user changes nothing.

\begin{longtable}{>{\raggedright\arraybackslash}p{0.4600\linewidth}>{\raggedright\arraybackslash}p{0.4600\linewidth}}
System operation & Refinement in the final design \\
\hline\endfirsthead
System operation & Refinement in the final design \\
\hline\endhead
\texttt{takeDelivery} & Taking a delivery is a guarded UPDATE: only one of two competing shippers succeeds. \\
\texttt{confirmDelivered} & Only the assigned shipper can finish; COD cash must equal the order total. \\
\texttt{reportFailure} & A failed delivery voids COD or records the full online amount as refund owed. \\
\end{longtable}
Editable sources: \texttt{models/\allowbreak{}UC-02/\allowbreak{}dcd\_\allowbreak{}uc02\_\allowbreak{}final\_\allowbreak{}overview.\allowbreak{}dot}, \texttt{models/\allowbreak{}UC-02/\allowbreak{}dcd\_\allowbreak{}uc02\_\allowbreak{}final\_\allowbreak{}detail\_\allowbreak{}A1.\allowbreak{}dot}, \texttt{models/\allowbreak{}UC-02/\allowbreak{}dcd\_\allowbreak{}uc02\_\allowbreak{}final\_\allowbreak{}detail\_\allowbreak{}A2.\allowbreak{}dot}, \texttt{models/\allowbreak{}UC-02/\allowbreak{}dcd\_\allowbreak{}uc02\_\allowbreak{}final\_\allowbreak{}detail\_\allowbreak{}A3.\allowbreak{}dot} and \texttt{models/\allowbreak{}UC-02/\allowbreak{}dcd\_\allowbreak{}uc02\_\allowbreak{}final\_\allowbreak{}detail\_\allowbreak{}A4.\allowbreak{}dot} (Graphviz, generated by \texttt{models/\allowbreak{}shared/\allowbreak{}final\_\allowbreak{}dcd.\allowbreak{}py}), with SVG exports beside them.

\textbf{Integration repair evidence}

The diagrams and typed metadata were refreshed against the corrected shared model. Allowed layered dependencies are drawn; the twelve retained dependency-policy entries are documented as static policy checks in report 4.2 and Table P4-1, not as runtime tests. P2-S13/P3-S03 identified C0 targets for guarded writes, commit failures and busy handling; their final scoped evidence is recorded in Table P4-1. Read paths remain read-only.

DeliveryController calls the Delivery and Payment domain receivers. The local variable delivery is resolved to Delivery; it is not an upward call to the delivery blueprint.

Diagram notation (R1--R4): every displayed association now has separate multiplicity text at its actual attached source and target ends, plus a separate role label. The initial SSD request precedes actual accepted/refused interaction operands; both success and refusal response paths are visible. The controller refusal reply is conditional on controller invocation. Read guards state no writes; command success requires commit. Call heads are closed/filled and reply heads are dashed/open.
